
De $s\ge1$, r\'esulte $G'\ne0$\,; par suite, $F'$ n'est pas contenu dans $F(-1)$\,; c'est-\`a-dire, la section $\sigma$ n'est pas contenue dans l'image de l'injection $H^0(F(-1))\to H^0(F)$, d'o\`u (i).

L'assertion (ii) r\'esulte imm\'ediatement de (i), compte tenu de ce que $h^0(F)=0$ implique \'evidemment $h^0(F(-1))=0$, car $h^0(F(-1))\le h^0(F)$. Dans (iii), on trouve $0\le h^0(F(n))-h^0(F(n-1))\le h^0(G(n))$. Donc pour $n\ge n_0$, $h^0(F(-n_0))=h^0(F(-n))$ car $h^0(G(-n_0))=0$\,; donc $h^0(F(-n_0))=0$ d'apr\`es (ii). Par suite, $h^0(F(p-n_0))\le h^0(G(1-n_0))+\ldots+h^0(G(p-n_0))\le ph^0(G(p-n_0))$ pour $p\ge1$\,; d'o\`u (iii).

\textbf{Proposition 6.6.} \emph{Sous les conditions de 6.1, on suppose $F$ un $(b_0,\ldots,b_r)$-faisceau. Alors $F_q$ est un $\bigl((b_{q-1})^q,\ldots,(b_{q-1})^2,b_{q-1},b_q,\ldots,b_r\bigr)$-faisceau.}

En effet, soit $\sigma$ une section ``g\'en\'erale'' de ${O}_P(1)$ ; alors $\sigma$ d\'efinit un diagramme commutatif \`a lignes et colonnes exactes
\[
\begin{tikzcd}[column sep=large,row sep=large]
&0 \arrow[d]&0 \arrow[d]&0 \arrow[d]&\\
0 \arrow[r]&N_q(-1) \arrow[r] \arrow[d]&N_q \arrow[r] \arrow[d]&G' \arrow[r] \arrow[d]&0\\
0 \arrow[r]&F(-1) \arrow[r] \arrow[d]&F \arrow[r] \arrow[d]&G \arrow[r] \arrow[d]&0\\
0 \arrow[r]&F_q(-1) \arrow[r] \arrow[d]&F_q \arrow[r] \arrow[d]&G'' \arrow[r] \arrow[d]&0\\
&0&0&0&
\end{tikzcd}
\]
o\`u $G$ est un $(b_1,\ldots,b_r)$-faisceau (1.6 (i)) et o\`u $G''=G_{q-1}$ (question de profondeur).
